% ============================================================
% PART VII: SYSTEM DESIGN (BIG INTEGRATED QUESTION)
%
% PEANUT EXAM DESIGN QUESTION LAYOUT STANDARD
%
% AI INSTRUCTION:
% - ONE large scenario that integrates most topics of the course
% - The background sits in \noindent\fbox{\parbox{\textwidth}{...}}
%   with an itemized list of hard requirements (numbers, latencies, scale)
% - Sub-questions are (a), (b), (c), ... each targeting a DIFFERENT topic,
%   with a bold title line, then the specific things to address
% - Requirements must be concrete and quantified so answers can be graded
% - \vspace after each sub-question is large (80-90mm): this is where
%   the student designs
% - Close with a trade-off / "what would you cut" question
% ============================================================

\section*{Part VII: System Design}

\subsection*{\textit{Title of the System to Be Designed}}

\vspace{3mm}

% The \dimexpr keeps the box inside the margins: an \fbox adds
% 2\fboxsep + 2\fboxrule on top of the parbox width.
\noindent\fbox{%
\parbox{\dimexpr\textwidth-2\fboxsep-2\fboxrule\relax}{%
\textbf{Background:}
One or two sentences naming the organization and the goal.
The system must:

\begin{itemize}
    \item Handle a \textbf{quantified ingest or workload}
    (e.g. 10,000+ events/sec from 10,000 devices).
    \item Meet a \textbf{latency or availability target}
    (e.g. detect within $\leq 50\,\text{ms}$, respond within $\leq 100\,\text{ms}$).
    \item Support a \textbf{batch or analytical path} alongside the real-time path.
    \item Serve \textbf{multiple tenants or departments} with isolation guarantees.
    \item Satisfy a \textbf{security or privacy constraint} that rules out the
    naive solution.
\end{itemize}
}}

\vspace{6mm}

Answer the following design sub-questions.
\textbf{Your answers must reference specific technologies and concepts from this course.}
Justify every major design decision.

\vspace{5mm}

%---------------------------------------------------------------
\begin{enumerate}[label=(\alph*)]

\item \textbf{First Layer of the Design} \\
Design this layer. Specifically:
\begin{itemize}
    \item Identify at least \textbf{three concrete components} and justify each.
    \item Explain how two specific course concepts apply here.
    \item State the guarantee you chose and the trade-off it costs you.
\end{itemize}

\vspace{85mm}

%---------------------------------------------------------------
\newpage
\item \textbf{Second Layer of the Design} \\
Address the requirement that the naive solution cannot satisfy.
Your answer must include a diagram or a labelled data flow.

\vspace{90mm}

%---------------------------------------------------------------
\item \textbf{Third Layer of the Design} \\
Handle the multi-tenant / isolation / scheduling requirement.
State what happens under contention.

\vspace{85mm}

%---------------------------------------------------------------
\newpage
\item \textbf{Failure and Trade-offs} \\
One component fails. Describe the blast radius, the recovery path,
and which requirement you would relax first if the budget were halved.

\vspace{90mm}

\end{enumerate}
